\begin{lemma}
Let $\sigma=(\sigma^1,\dots,\sigma^m)$ be vector fields $\sigma^i:\mathbb{R}^n\to\mathbb{R}^n$, let
$h\in C^2(\mathbb{R})$, let $w:\mathbb{R}^n\to\mathbb{R}$ and let $x\in\mathbb{R}^n$. Assume that
$w$ is of class $C^2$ in a neighbourhood of $x$ and that every $\sigma^i$ is of class $C^1$ in a
neighbourhood of $x$. Then
$$\mathfrak{X}(h\circ w)(x)=h'\big(w(x)\big)\,\mathfrak{X}w(x)$$
and
$$\mathfrak{X}\mathfrak{X}(h\circ w)(x)=h'\big(w(x)\big)\,\mathfrak{X}\mathfrak{X}w(x)
+h''\big(w(x)\big)\,\mathfrak{X}w(x)\otimes\mathfrak{X}w(x),$$
where $\xi\otimes\eta\in\mathbb{R}^{m\times m}$ denotes the matrix with entries
$(\xi\otimes\eta)_{i,j}=\xi_i\,\eta_j$. These are the two identities
$\mathfrak{X}(h\circ w)=h'(w)\mathfrak{X}w$ and $\mathfrak{X}\mathfrak{X}(h\circ
w)=h'(w)\mathfrak{X}\mathfrak{X}w+h''(w)\,\mathfrak{X}w\otimes\mathfrak{X}w$ with which the proof
of Lemma \texttt{lem:tech} of the paper begins.
\end{lemma}
\begin{proof}
Write $t:=h'(w(x))$, $s:=h''(w(x))$ and $\xi:=\mathfrak{X}w(x)$.

\emph{Step 1: the first identity.} Let $V$ be an open neighbourhood of $x$ on which $w$ is of class
$C^2$. Since $h$ is differentiable on $\mathbb{R}$ and $w$ is differentiable on $V$, the chain
rule gives that $h\circ w$ is differentiable on $V$ with
\begin{equation*}D(h\circ w)(y)=h'\big(w(y)\big)\,Dw(y)\qquad\text{for every }y\in V. \tag{1}\end{equation*}
Evaluating (1) at $y=x$ and at the vector $\sigma^i(x)$, and using the definition of
\noderef{HGrad},
$$\big(\mathfrak{X}(h\circ w)(x)\big)_i=D(h\circ w)(x)\big[\sigma^i(x)\big]
=t\,Dw(x)\big[\sigma^i(x)\big]=t\,\xi_i,$$
for every $i$, which is the first identity.

\emph{Step 2: the second derivative of $h\circ w$.} Since $h\in C^2(\mathbb{R})$ and $w$ is of
class $C^2$ on $V$, the composition $h\circ w$ is of class $C^2$ on $V$, hence in a neighbourhood
of $x$. Moreover $h\in C^2(\mathbb{R})$ and $w$ is of class $C^2$ in a neighbourhood of $x$, which
are the hypotheses of \noderef{SndFDerivCompScalar}; applied to $h$, $w$ and $x$ it gives
\begin{equation*}D^2(h\circ w)(x)[z_1,z_2]=s\,Dw(x)[z_1]\,Dw(x)[z_2]+t\,D^2w(x)[z_1,z_2] \tag{2}\end{equation*}
for all $z_1,z_2\in\mathbb{R}^n$.

\emph{Step 3: the second identity.} Fix $i,j$. By Step 2 the function $h\circ w$ is of class $C^2$
in a neighbourhood of $x$, and $\sigma^i$ is of class $C^1$ in a neighbourhood of $x$; so
\noderef{HHessEntryFormula} applies to $h\circ w$ and gives
$$\big(\mathfrak{X}\mathfrak{X}(h\circ w)(x)\big)_{i,j}
=D^2(h\circ w)(x)\big[\sigma^j(x),\sigma^i(x)\big]
+D(h\circ w)(x)\big[D\sigma^i(x)\,\sigma^j(x)\big].$$
By (2) with $z_1=\sigma^j(x)$, $z_2=\sigma^i(x)$, and by the definition of \noderef{HGrad}, the
first summand equals
$$s\,Dw(x)\big[\sigma^j(x)\big]\,Dw(x)\big[\sigma^i(x)\big]
+t\,D^2w(x)\big[\sigma^j(x),\sigma^i(x)\big]
=s\,\xi_j\,\xi_i+t\,D^2w(x)\big[\sigma^j(x),\sigma^i(x)\big],$$
while by (1) the second summand equals
$t\,Dw(x)\big[D\sigma^i(x)\,\sigma^j(x)\big]$. Adding and applying
\noderef{HHessEntryFormula} to $w$ itself (which is of class $C^2$ near $x$),
$$\big(\mathfrak{X}\mathfrak{X}(h\circ w)(x)\big)_{i,j}
=t\Big(D^2w(x)\big[\sigma^j(x),\sigma^i(x)\big]+Dw(x)\big[D\sigma^i(x)\,\sigma^j(x)\big]\Big)
+s\,\xi_i\,\xi_j
=t\,\big(\mathfrak{X}\mathfrak{X}w(x)\big)_{i,j}+s\,\xi_i\,\xi_j .$$
Since $(\xi\otimes\xi)_{i,j}=\xi_i\xi_j$, this is precisely the second identity entrywise, and
hence as an identity of matrices.
\end{proof}
